


\chapter{Revisão de Literatura}\label{chap:levantamento_literatura}

Este capítulo posiciona a proposta frente a duas frentes de literatura: os sistemas de revisão científica assistida por LLMs e as ferramentas de extração estruturada de documentos científicos.

\section{Revisão científica assistida por agentes e LLMs}

Investigações recentes exploram o papel de agentes baseados em LLMs no processo de revisão por pares. O framework AgentReview utiliza agentes para simular as interações entre autores, revisores e editores, permitindo investigar fatores sociológicos e organizacionais que influenciam decisões editoriais, e evidenciando a viabilidade de arquiteturas multiagentes para tarefas avaliativas \cite{AgentReview}. Na geração automática de pareceres, o método MARG (Multi-Agent Review Generation) propõe múltiplos agentes especializados que analisam diferentes aspectos do artigo, como experimentos, clareza e impacto \cite{MARG}.

Avanços posteriores aprofundam a qualidade das revisões por meio de raciocínio estruturado em múltiplos estágios. O framework DeepReview emula o processo de revisores especialistas, com verificação de novidade, avaliação multidimensional e verificação de confiabilidade, enfatizando que o objetivo não é substituir o revisor humano, mas explorar o potencial de assistência das LLMs \cite{DeepReview}. Em contrapartida, estudos empíricos revelam vulnerabilidades desses sistemas: \cite{WhereLLMWrong} propõem um framework de perturbação multinível guiado por aspectos para diagnosticar sistematicamente as fraquezas de LLMs na revisão automatizada.

Em conjunto, esses trabalhos sugerem que agentes de IA podem desempenhar papel relevante no apoio à produção científica, mas evidenciam limitações que precisam ser consideradas no projeto de arquiteturas multiagentes. Persistem lacunas quanto à adaptação dessas tecnologias a contextos editoriais específicos, sobretudo fora do eixo das grandes conferências internacionais. A presente proposta distingue-se dessa literatura em dois aspectos: primeiro, ao deslocar o objetivo de avaliar a qualidade absoluta para verificar a conformidade a um gênero-alvo, o que ataca diretamente a rejeição precoce; segundo, ao fundamentar essa verificação em uma modelagem explícita e auditável do gênero da trilha, e não apenas no raciocínio implícito de um modelo.

\section{Extração estruturada de documentos científicos}

O método proposto depende de converter cada PDF em uma representação estruturada antes de qualquer raciocínio por LLM. Nesse domínio, duas famílias de ferramentas recentes voltadas a pipelines de IA generativa, como o Docling \cite{Docling}, toolkit de código aberto que converte PDF e outros formatos em representação estruturada unificada, com modelos especializados de análise de layout e reconhecimento de tabelas, e que preserva figuras e tabelas. 
%Como se detalha no Capítulo 5, 
Este trabalho adota o Docling, por operar inteiramente em ambiente Python, preservar elementos visuais relevantes para a caracterização de gênero e oferecer melhor suporte ao pós-processamento de textos em português, complementando-o com consultas às APIs CrossRef e Semantic Scholar para o tratamento de referências.

Para a recuperação semântica, adota-se o modelo de \textit{embeddings} BGE-M3 \cite{BGEM3}, multilíngue, com forte desempenho em português, capaz de recuperação densa e esparsa e de processar contextos de até 8.192 tokens — características adequadas a um corpus em língua portuguesa e passíveis de execução local.